\documentclass[a4paper,11pt]{article}

\usepackage[ngerman]{babel}
\usepackage[top=3cm,bottom=3cm,left=3cm,right=3cm]{geometry} %NICHT IN KEMIE2


%\usepackage{microtype} %schöneres typesetting  mit [stretch=10] für nur 1% anstelle von 2%
%\usepackage{lmodern} %Latin Modern FONT
%\usepackage{pifont} %schönere Font



\usepackage{amsmath}
\usepackage{nicefrac}
\usepackage{amssymb}
%\usepackage{mathtools}

\usepackage{titlesec} %Um Section, etc bearbeiten zu können
\usepackage{setspace}
\usepackage{enumitem} %enumerate
\usepackage{multicol}
\usepackage[many]{tcolorbox}
\usepackage{hyperref}
\usepackage{arydshln}

%\renewcommand{\ttdefault}{lmtt} %Schriftart wird geändert zu Latin Modern Typewriter


\hypersetup{hidelinks,
			colorlinks=true,
			linkcolor=black,
			citecolor=miudiutex,
			urlcolor=miugray2
			}

\setlength{\parindent}{0pt}
\setcounter{section}{0}


\titlespacing*{\section}{0pt}{20mm}{6mm}
\titlespacing*{\subsection}{0pt}{6mm}{3mm}
\titlespacing*{\subsubsection}{0pt}{3mm}{0mm}



%\definecolor{miudiutex}{HTML}{2f3d39}
\definecolor{miudiutex}{HTML}{1d2926}
\definecolor{miugray}{HTML}{b4cccf}
\definecolor{miugray2}{HTML}{4d7365}
\definecolor{red}{HTML}{cf1f5c}
\definecolor{green}{HTML}{00A36C}
\definecolor{delphinidin}{HTML}{315794}
\definecolor{pelargonidin}{HTML}{b33807}
\definecolor{cyanidin}{HTML}{ba0746}
\definecolor{petunidin}{HTML}{b56fed}



\raggedcolumns %wichtig für Multicolumns


%================================
% SYMBOLS
%================================

\newcommand{\RA}{$\rightarrow$~}
\newcommand{\RL}{$\rightleftharpoons$~}
\newcommand{\HARP}{\ding{226}~}
%$\xrightarrow{2}$ %Pfeil mit 2 drüber


%================================
% SHORT COMMANDS (BOXES ; ETC)
%================================


\newcommand{\notiz}[1]{\begin{notizz}{}{}\begin{center}
			{\color{miugray2!50!miudiutex}#1} \end{center}\end{notizz}}
\newcommand{\zit}[2]{\begin{zitat*}{#1}{}\small #2 \end{zitat*}}

\newcommand{\frage}[2]{\begin{qst*}{#1}{}\small #2 \end{qst*}}

\newcommand{\unten}{\end{center} \tcblower \begin{center}}


\newcommand*\circled[1]{\tikz[baseline=(char.base)]{
            \node[shape=circle,draw,inner sep=0.5pt,miudiutex] (char) {\tiny #1};}}
            
\newcommand{\miusquare}{{\color{miugray2!50}\scriptsize $\blacksquare$}~}
\newcommand{\miusquarp}{{\color{petunidin!50}\scriptsize $\blacktriangleright$}~}
\newcommand{\niu}{\item[\miusquare]}
\newcommand{\nip}{\item[\miusquarp]}

%%%%% ABSAETZE

\newcommand{\pps}{\par \vspace*{1mm}}
\newcommand{\pp}{\par \vspace*{3mm}}
\newcommand{\ppbig}{\par \vspace*{8mm}}
\newcommand{\pphuge}{\par \vspace*{10mm}}

%================================
% FRAGE BOX
%================================

\tcbuselibrary{theorems,skins,hooks}
\newtcbtheorem{qst}{Frage:}
{%
	enhanced,
	breakable,
	colback = gray!15!white,
	frame hidden,
	boxrule = 0sp,
	borderline west = {2.5pt}{0pt}{red!70},
	sharp corners,
	detach title,
	before upper = \tcbtitle\par\smallskip,
	coltitle = black,
	fonttitle = \bfseries\sffamily,
	description font = \mdseries,
	separator sign none,
	segmentation style={solid, gray!50!black},
}
{th}




%================================
% ZITAT BOX
%================================

\tcbuselibrary{theorems,skins,hooks}
\newtcbtheorem{zitat}{Zitat}
{%
	enhanced,
	breakable,
	colback = gray!15!white,
	frame hidden,
	boxrule = 0sp,
	borderline west = {2.5pt}{0pt}{black!70},
	sharp corners,
	detach title,
%	before upper = \tcbtitle\par\smallskip,
	coltitle = gray!50!black,
	fonttitle = \bfseries\sffamily,
	description font = \mdseries,
%	separator sign none,
	segmentation style={solid, gray!50!black},
}
{th}



%================================
% SIMPLE SCHWARZER RAHMEN BOX
%================================


\newcommand{\boxfr}[1]{
	\begin{tcolorbox}[
	drop shadow southeast,
	enhanced,
	colframe=black!50,
	colback=white,
	boxrule = 1pt, 
%	toprule=0pt, bottomrule=0pt,rightrule=0pt,leftrule=0pt,
	boxsep=5pt,
	arc=1pt,
	]
	\begin{center}
	#1
	\end{center}
	\end{tcolorbox}
}

\definecolor{miudiutex}{HTML}{1d2926}
\definecolor{miugray}{HTML}{b4cccf}
\definecolor{miugray2}{HTML}{4d7365}
\definecolor{white2}{HTML}{f5f7f8}
\definecolor{red}{HTML}{cf1f5c}
\definecolor{green}{HTML}{00A36C}
\definecolor{delphinidin}{HTML}{315794}
\definecolor{uniblau}{HTML}{004e9f}
\definecolor{unigelb}{HTML}{fcba00}
\definecolor{unigrau}{HTML}{909085}
\definecolor{pelargonidin}{HTML}{b33807}
\definecolor{cyanidin}{HTML}{ba0746}
\definecolor{paeonidin}{HTML}{d15ca6}
\definecolor{petunidin}{HTML}{b56fed}
\definecolor{malvidin}{HTML}{8a2d8c}

\newcommand{\metermilli}{\mskip3mu \text{mm}}
\newcommand{\cm}{\mskip3mu \text{cm}}
\newcommand{\m}{ \text{m}}
\newcommand{\lite}{ \text{l}}
\newcommand{\kg}{ \text{kg}}
\newcommand{\g}{ \text{g}}
\newcommand{\km}{ \text{km}}
\newcommand{\h}{ \text{h}}
\newcommand{\s}{ \text{s}}
\newcommand{\tonne}{ \text{t}}
\usepackage[utf8]{inputenc}
\usepackage[T1]{fontenc}
\usepackage{lmodern}

\usepackage[
    activate={true,nocompatibility},
    final,
    tracking=true,
    kerning=true,
    spacing=true,
    factor=1100,
    stretch=30,
    shrink=20
]{microtype}
\usepackage{pdfpages}
\usepackage{awesomebox}
\usepackage{marginnote}
\tcbuselibrary{documentation}
\usepackage{sidenotes}
\usepackage{import}
\usepackage{xifthen}
\usepackage{transparent}
\usepackage[table]{xcolor}


\newcommand{\abbicon}{%
  \protect\tikz[baseline=0.2mm,scale=0.8]{%
    % Das blaue Quadrat mit abgerundeten Ecken
    \fill[uniblau, rounded corners=1.2pt] (0,0) rectangle (0.8em, 0.8em);
    % Der weiße Kreis in der Mitte
    \fill[white] (0.4em, 0.4em) circle (0.12em);
  }%
} 
\usepackage[labelfont={bf},font={color=miudiutex,small},figurename={\abbicon $~$Abbildung}]{caption}


\newcommand{\bckgrnd}{
\AddToHookNext{shipout/background}{
\begin{tikzpicture}[remember picture, overlay]
 \draw[draw=none,fill=uniblau!70, shift={(current page.north west)}] (-3,0) rectangle (23,-3.75);
 \draw[draw=none,fill=miugray, shift={(current page.north east)}] (-7,0) rectangle (3,-3.75);
\end{tikzpicture}
}
}



\newcommand{\incfig}[2]{%
\def\svgwidth{#2\columnwidth}
\import{C://LaTeX-Files/pngs/}{#1.pdf_tex}
}

\renewcommand{\textbf}[1]{\textsf{{\bfseries {#1}}}}



\titleformat{\section}[hang]
		{\LARGE \color{uniblau!75!black} \sffamily \bfseries}
		{\thesection}
		{6mm}
		{}
		[]
\titleformat{\subsection}[hang]
		{\large \sffamily \bfseries \color{black}}
		{\thesubsection}
		{4mm}
		{{\color{miugray}$\blacksquare ~$}}

\titleformat{\subsubsection}[block]
		{\normalsize \bfseries \sffamily \color{miudiutex}}%
		{\normalsize \color{black} \thesubsubsection}
		{2mm}
		{{\color{miugray}$\blacksquare ~$}}
		[ \color{black}]
		
\titleformat{\paragraph}
		{\normalsize \bfseries \sffamily \color{black}}
		{\footnotesize \theparagraph}
		{1mm}
		{}

\pagestyle{empty}
\titlespacing*{\section}{0pt}{20mm}{12mm}
\titlespacing{\section}{0pt}{20mm}{12mm}

\titlespacing*{\subsection}{0pt}{14mm}{4mm}
\titlespacing{\subsection}{0pt}{14mm}{4mm}

\titlespacing*{\subsubsection}{0pt}{6mm}{3mm}
\titlespacing{\subsubsection}{0pt}{6mm}{3mm}

\titlespacing*{\paragraph}{0pt}{6mm}{2mm}
\titlespacing{\paragraph}{0pt}{6mm}{2mm}


\let\oldmarginfigure\marginfigure
\let\endoldmarginfigure\endmarginfigure

\let\oldmarginfigure\marginfigure
\let\endoldmarginfigure\endmarginfigure

\renewenvironment{marginfigure}[1][0pt] % [1] steht für 1 Argument, [0pt] ist der Standardwert
  {\oldmarginfigure[#1]\captionsetup{labelfont={sf,bf,color=uniblau!75!black},font={color=miudiutex,small}}}
  {\endoldmarginfigure}
  
  
  
  
\renewcommand{\labelitemi}{\color{uniblau!70}$\bullet$}
\renewcommand{\labelitemii}{\color{uniblau!70}$\cdot$}
\newcommand{\plmtsquare}{{\color{uniblau!70}$\blacksquare~$}}
\newcommand{\splmt}[1]{{\color{uniblau!75!black} \sffamily #1}}
\newcommand{\fplmt}[1]{{\color{uniblau!75!black} \sffamily \bfseries #1}}

\newcommand{\antwort}[1]{
\begin{tcolorbox}[
	enhanced,
	enforce breakable,
	standard jigsaw,
	boxrule=0.7pt,
	colframe=black,
	sharp corners,
	boxsep=5pt,
	title=\bfseries \sffamily Antwort,
	opacityback=0,
	coltitle=miudiutex,
	attach boxed title to top text left={yshift=-\tcboxedtitleheight/2},
	boxed title style={colback=white,colframe=white},
	bottomrule at break=0pt,
	toprule at break=0pt,
	pad at break=16mm,
]
	\begin{center}
	#1
	\end{center}
	\end{tcolorbox}
}




\rowcolors{2}{miugray!50}{miugray!50}
\arrayrulecolor{white} % Weiße Gitterlinien
\setlength{\arrayrulewidth}{1.5pt} % Etwas dickere Linien
\renewcommand{\arraystretch}{1.4} % Mehr Zeilenhöhe für bessere Lesbarkeit}


%\setlist[itemize]{itemsep=0mm}
\setlist[enumerate]{itemsep=0mm}
\usepackage{relsize}
\usetikzlibrary{patterns, arrows.meta, calc}

\newif\ifshowme
%\showmetrue




\begin{document}
\newgeometry{top=1.8cm,bottom=4cm,left=1.5cm,right=7.5cm,marginparsep=-2.00cm,marginparwidth=4.75cm}


\bckgrnd


\begin{center}
\LARGE
\color{white}
\textbf{Übung 05} \par \vspace*{2mm}
\large \color{miugray} \textbf{Prozessbezogene Lebensmitteltechnologie}
\end{center}
\par 
\vspace*{6mm}


\color{miudiutex}

\section*{Wärmeübertragung II}
\subsection*{Aufgabe 1}
Ein 2 cm dickes Stahlrohr (Wärmeleitfähigkeit = 43 W/(m°C)) mit einem
Innendurchmesser von 6 cm wird verwendet, um Dampf von einem Kessel über eine
Strecke von 40 m zu einer Prozessanlage zu leiten. Die Oberflächentemperatur der
Innenseite des Rohrs beträgt 115 °C und die Oberflächentemperatur der Außenseite
90 °C. Berechnen Sie den gesamten Wärmeverlust an die Umgebung unter stationären
Bedingungen.
\begin{marginfigure}[-3cm]
\centering
\resizebox{.6\textwidth}{!}{
\includegraphics[width=12cm]{F://LaTeX-Files/pngs/40mrohr}
}
\caption{Modell eines Stahlrohres}
\end{marginfigure}


\ifshowme
\antwort{
\begin{align*}
\dot{Q} &= (T_i - T_a) \frac{2 \pi l \cdot \lambda}{\ln(\frac{r_a}{r_i})} \\
		&= (115 - 90)\text{°C} \cdot \frac{2 \pi \cdot 40~\text{m} \cdot 43~ \text{W/m°C}}{\ln(\frac{0,05~\text{m}}{0,03~\text{m}})} \\
		&= 528.902,54~ \text{W}
\end{align*}
}
\pagebreak
\else
\fi



\subsection*{Aufgabe 2}
Eine Kühlhauswand (3 m x 6 m) wird aus 15 cm dickem Beton (Wärmeleitfähigkeit = 1,37
W/(m°C)) erstellt. Es soll auch eine Isolierung an der Innenseite eingebaut werden, um
die Wärmeübertragungsrate durch die Wand bei oder unter 500 W zu halten. Berechnen
Sie die erforderliche Dicke der Dämmung. Die Wärmeleitfähigkeit der Dämmung liegt bei
0,04 W/(m°C). Nehmen Sie an, dass die Temperatur an der Außenseite 38 °C und die
Temperatur an der Oberfläche der Dämmung 5 °C beträgt.
\begin{marginfigure}[-3.5cm]
\centering
\resizebox{\textwidth}{!}{
\begin{tikzpicture}

    % Schicht 1 (schraffiert) - Breite x
    \draw[thick,fill=miugray!50] (0,0) rectangle (1,5);

    % Schicht 2 - Breite 15cm
    \draw[thick] (1,0) rectangle (3,5);

    % Wärmestrom Pfeil (500 W)
    \draw[very thick, -{Stealth[scale=1.3]}] (3.5,3) -- (-1,3) node[above,yshift=2mm,
 font=\large
    ] {500 W};

    % Temperaturen
    \node[font=\large] at (-0.8,2) {5°C};
    \node[font=\large] at (3.8,2) {38°C};

    % Bemaßung unten (x)
    \draw[<->,thick, >=latex] (0,-0.3) -- (1,-0.3) node[midway, below,font=\large] {$x$};
    
    % Bemaßung unten (15cm)
    \draw[<->,thick, >=latex] (1,-0.3) -- (3,-0.3) node[midway, below,font=\large] {15 cm};

\end{tikzpicture}
}
\caption{Modell einer Kühlhauswand mit Dämmung}
\end{marginfigure}

\ifshowme%
\marginnote[]{\small Hier ist die Verwendung der richtigen Vorzeichen besonders wichtig. Hätten wir für die Temperatur einen negativen Wert raus, müssten wir auch einen negativen Vorfaktor vor dem Bruch benutzen. Ein falsches Vorzeichen würde tatsächlich auch eine andere Zahl im Ergebnis zur Folge haben.}[2cm]%
%
\antwort{
\begin{align*}
\dot{Q} &= \frac{A \cdot T_a - T_i}{\frac{d_2}{\lambda_2} + \frac{d_1}{\lambda_1}} \\
	d_2 &= \lambda_2 \cdot (\frac{A \cdot (T_a - T_i)}{\dot{Q}} - \frac{d_1}{\lambda_1}) \\
	d_2 &= 0,04~ \text{W/m°C} \cdot (\frac{18~\text{m}^2 \cdot ((38-5) \text{°C})}{500~\text{W}} - \frac{0,15~\text{m}}{1,37~ \text{W/m°C} }) \\
	d_2 &= 0,043~\text{m} \rightarrow 4,3~ \text{cm}
\end{align*}
}
\else
\fi

\subsection*{Aufgabe 3}

\begin{enumerate}[label=\alph*),itemsep=-3mm]

\item 
Wie groß ist der Wärmestrom durch ein 5 m langes Rohr aus Stahl mit
Innendurchmesser 10 cm und Außendurchmesser 11 cm, wenn im Innern eine
Flüssigkeit der Temperatur 80 °C und außen eine Flüssigkeit der Temperatur
20 °C zirkuliert? Stahl hat eine Wärmeleitfähigkeit von 16,3 W/(mK). Die
Widerstände der Wärmeübergänge an der inneren und äußeren Rohroberfläche
seien zu vernachlässigen.




\begin{marginfigure}[-3cm]
\begin{tikzpicture}[scale=0.6]
\draw[line width=0.01mm,fill=miugray] (0,0) circle (2cm);
\draw[miugray2,fill=miugray!50] (0,0) circle (1.5cm);
\draw[miugray2,fill=white] (0,0) circle (1.3cm);
\draw (0,0) to (20:2cm) node [right] {$r_\mathsmaller{{SA}}$ = 0,055 m};
\draw[miugray!75!black] (0,0) to (-20:1.5cm) node [right,xshift=4mm,yshift=0mm] {$r_\mathsmaller{{SI}}$ = 0,05 m};
\draw[uniblau] (0,0) to (-60:1.3cm) node [right,xshift=4mm,yshift=-4mm] {$r_\mathsmaller{{GI}}$ = 0,047 m};
\end{tikzpicture}
\caption{Stahlrohr mit innerer Glasschicht.}
\end{marginfigure}


\ifshowme
\antwort{
\begin{align*}
\dot{Q} &= (T_i - T_a) \frac{2 \pi l \cdot \lambda}{\ln(\frac{r_i}{r_a})} \\
\dot{Q} &= 60~ \text{°C} \cdot \frac{2 \pi \cdot 5~\text{m} \cdot 16,3~ \text{W/m°C} }{\ln(\frac{55}{50})} \\
	&= 322.366,15~ \text{W}
\end{align*}
}



\pagebreak
\else
\fi


\item Wie ändert sich der Wärmestrom, wenn innen eine 0,3 cm dicke Schicht aus
Glas aufgetragen wird? Die Wärmeleitfähigkeit von Glas liegt bei 1,0 W/(mK).
\end{enumerate}

\ifshowme
\antwort{
\begin{align*}
\dot{Q} &= 2 \pi \cdot l \cdot (T_i - T_a) \cdot \left( \frac{\ln(\frac{r_{SA}}{r_{SI}})}{\lambda_1} + \frac{\ln(\frac{r_{GA}}{r_{GI}})}{\lambda_2} \right)^{-1} \\
\dot{Q} &= 2 \pi \cdot 5m \cdot 60~ \text{°C} \cdot \left( {\frac{\ln(\frac{0,055}{0,05})}{16,3~ \text{W/m°C}} + \frac{\ln(\frac{0,05}{0,047})}{1~ \text{W/m°C}}} \right)^{-1} \\
	&= 27.833,46~ \text{W}
\end{align*}
}
\else
\fi




\subsection*{Aufgabe 4}
\begin{enumerate}[label=\alph*)]
\item Wie groß ist der Wärmeübergangskoeffizient innen, wenn Wasser mit einer
Geschwindigkeit von 0,5 m/s durch das Stahlrohr von Aufgabe 3 fließt?

\begin{table}[h]
\centering
\caption{Eigenschaften von Wasser bei 80 °C}
%\resizebox{\textwidth}{!}{
\begin{tabular}{l|l}\rowcolor{miugray!50}
Dichte & $\rho$ = 972kg/m$^3$ \\ \hline
Kinematische Viskosität & $\nu$  = 0,3648 x 10$^{-6}$ m$^{2}$/s \\ \hline
Wärmeleitfähigkeit & $\lambda$ = 0,67 W/(mK) \\ \hline
Prandtlzahl & Pr = 2,22  \\ \hline
\end{tabular}
%}
\end{table}

\ifshowme
\antwort{
Gesucht ist $\alpha$
\begin{align*}
\text{Nu} = \frac{\alpha \cdot d}{\lambda}
\end{align*}
\begin{align} \label{alpha}
\alpha = \frac{\text{Nu} \cdot \lambda}{d}
\end{align}
Für die Berechnung von $\alpha$ fehlt uns die Nusseltzahl Nu:
\begin{align}\label{nus}
\text{Nu} = \frac{(\frac{\xi}{8}) \cdot \text{Re} \cdot \text{Pr}}{1 + 12,7 \cdot (\frac{\xi}{8})^{\frac{1}{2}} \cdot (\text{Pr}^{\frac{2}{3}}-1)} \cdot \left[ 1 + (\frac{d}{l})^{\frac{2}{3}} \right]
\end{align}
Wobei
\begin{align*}
\xi &= (1,8 \cdot \log_{10} (\text{Re}) - 1,5)^{-2}
\end{align*}
Und Reynoldszahl Re:
\marginnote[]{Hier ist mit $l$ nicht die Länge des Rohres, sondern die charakteristische Länge des geometrischen Objektes gemeint. Dies ist beim Rohr i.d.R der Innendurchmesser}[1cm]
\begin{align*}
\text{Re} &= \frac{v \cdot l}{\nu} \\
\text{Re} &= \frac{0,5~ \text{m/s} \cdot 0,1~\text{m}}{0,3648 \cdot 10^{-6} ~\text{m}^2 /s} \\
&= 137.061,4035
\end{align*}
Daraus folgt für $\xi$:
\begin{align*}
\xi &= (1,8 \cdot \log_{10} (137.061) - 1,5)^{-2} \\
 &= 0,0167
\end{align*}
Werte einsetzen in (\ref{nus}):
\marginnote[]{hier ist d = Rohrinnendurchmesser und l = Länge des Rohres}[2cm]
\begin{align*}
\text{Nu} &= \frac{(\frac{0,0167}{8}) \cdot  137.061 \cdot 2,22}{1 + 12,7 \cdot \sqrt{\frac{0,0167}{8}} \cdot (2,22^{\frac{2}{3}}-1)} \cdot \left[ 1 + (\frac{0,094 ~ \text{m}}{5 ~ \text{m}})^{\frac{2}{3}} \right] \\
 &= 451,373 \cdot 1,070703118 \\
 &\approx 483
\end{align*}
Daraus folgt:
\begin{align*}
\alpha &= \frac{483 \cdot 0,67 \text{W/(mK)}}{ 0,1~ \text{m}} \\
 &\approx 3.236,1 ~ \text{W/(m$^2$ K)}
\end{align*}
}
\else
\fi

\ifshowme
\pagebreak
\newgeometry{top=1cm,bottom=1cm,left=1cm,right=4cm}
\else
\fi

\item Wie ändert sich der Wärmestrom der Aufgabe 3 unter Berücksichtigung des
Widerstands des inneren Wärmeübergangs? \pp

\ifshowme
\antwort{
Wärmedurchgang - Rohr:
\begin{align*}
\dot{Q} &= 2 \pi \cdot l \cdot \Delta T \cdot \left( \frac{1}{\alpha_i \cdot r_i} + \frac{1}{\alpha_a \cdot r_a} + \frac{\ln(\frac{r_a}{r_i})}{\lambda_W} \right)^{-1} \\
\end{align*}
Hier stellt sich aber die Problematik, dass wir $\alpha_a$ nicht kennen und es deswegen vernachlässigen müssen. Desweiteren ist $\lambda_W$ die Summe aus $\lambda_G$ für die innere Glasschicht und $\lambda_S$ für die äußere Stahlschicht. \par 
Daraus ergibt sich:
\begin{align*}
\dot{Q} &= 2 \pi \cdot l \cdot \Delta T \cdot \left( \frac{1}{\alpha_i \cdot r_i} +  \frac{\ln(\frac{r_{aG}}{r_{iG}})}{\lambda_G} +  \frac{\ln(\frac{r_{aS}}{r_{iS}})}{\lambda_S} \right)^{-1}
\end{align*}
Wenn wir nun die Werte einsetzen: 
\begin{align*}
\dot{Q} &= 2 \pi \cdot 5~ \text{m} \cdot 60\text{°C} \cdot \left( \frac{1}{3.266,95 ~ \text{W/(m$^2$ K)} \cdot 0,047 ~ \text{m}} +  \frac{\ln(\frac{0,055}{0,05})}{16,3~ \text{W/m°C}} + \frac{\ln(\frac{0,05}{0,047})}{1~ \text{W/m°C}} \right)^{-1} \\
 &= 25.391,62~ \text{W}
\end{align*}
}
\else
\fi

\ifshowme

\antwort{
\begin{align*}
\dot{Q} &= 2 \pi \cdot l \cdot \Delta T \cdot \left( \frac{1}{\alpha_i \cdot r_i} + \frac{1}{\alpha_a \cdot r_a} + \frac{\ln(\frac{r_a}{r_i})}{\lambda_W} \right)^{-1} \\
\end{align*}
Hier stellt sich aber die Problematik, dass wir $\alpha_a$ nicht kennen und es deswegen vernachlässigen müssen. \par 
Daraus ergibt sich:
\begin{align*}
\dot{Q} &= 2 \pi \cdot l \cdot \Delta T \cdot \left( \frac{1}{\alpha_i \cdot r_i} +   \frac{\ln(\frac{r_{aS}}{r_{iS}})}{\lambda_S} \right)^{-1}
\end{align*}
Wenn wir nun die Werte einsetzen: 
\begin{align*}
\dot{Q} &=  \frac{\pi \cdot 5~ \text{m} \cdot 60\text{°C} }{\frac{1}{3.245,28 ~ \text{W/(m$^2$ K)} \cdot 0,05 ~ \text{m}} +  \frac{\ln(\frac{0,055}{0,05})}{16,3~ \text{W/m°C}}}    \\
 &= 156.948,23~ \text{W} \\
 &= 157 ~\text{kW}
\end{align*}
}
\else
\fi

\end{enumerate}






\end{document}
